\section{Implications}

If the model holds, several deep things follow. They are not decorations on the edges of the theory. They are the payoff. The base is a single substance, the vibe, valued in each site by one tone. From that one substance, time, space, expansion, unification, and feeling all turn out to be the same growth read in different ways. This section collects those consequences and states plainly what is at stake if the model is right.

\begin{principle}[The one substance, read five ways]
The base is one substance, the vibe, carried by the eight atoms (mesh, dock, site, tone, lean, knit, wake, beat). From it, time, space, cosmic expansion, force unification, and feeling are not five separate facts. They are one growth read five different ways.
\end{principle}

Two terms recur throughout, so a plain statement of each helps before we begin. A \emph{beat} is the discrete step of time, the single tick at which the knit runs once and the wake grows once. The \emph{wake} is the growing edge of the mesh, the open frontier where new docks come into being at peace. These two atoms, the when and the how-the-space-changes, carry most of the weight in what follows.

\subsection{Time, the arrow, and expansion are one process}

The mesh grows one shell of docks per beat. That single act of growth is three things at once, depending on which feature of it you attend to. It is not that the model bundles three mechanisms together. There is one mechanism, the growth at the wake, and three names for three readings of it.

\begin{table}[ht]
\centering
\begin{tabular}{@{}ll@{}}
\toprule
reading & what it is \\
\midrule
the beat & one ring of growth is one tick of time \\
the arrow & the growth runs one way, and that direction is the arrow \\
the expansion & the spatial volume per beat grows, and that is cosmic expansion \\
\bottomrule
\end{tabular}
\caption{Three readings of one growth at the wake.}
\label{tab:growth-readings}
\end{table}

Read the first way, each fresh ring of docks is one tick. The beat is the growth counted. Read the second way, the growth runs outward, from the trunk toward the open frontier, and never the other way. That one-way running is the arrow of time, emergent from the wake acting through the lean, not a direction built into the base. The bulk knit is reversible and carries no preferred direction of its own. Read the third way, the spatial volume added per beat keeps growing, because the mesh is hyperbolic and its room grows exponentially. That growth of volume per beat is cosmic expansion.

\begin{result}[Time is the growing]
Cosmic expansion is not added on top of the mesh. It \emph{is} the mesh growing at its wake. Time \emph{is} that growing, counted in beats. The arrow is the one-way sense of the same growth. Three usually separate facts collapse into one process.
\end{result}

Three facts that are normally posited and tuned independently, the flow of time, the direction of time, and the expansion of space, here come out of a single growth. There is nothing to tune among them, because there is only one thing happening.

\subsection{An age and a horizon, on either reading}

The model does not force a choice between a finite and an infinite past. Both readings of birth are live, and the payoff lands either way. This is the open question told in full elsewhere in the paper, but the implication for age and horizon can be stated cleanly here.

On the finite reading the mesh has a true first beat and a finite size. There is a first ring of docks, and counting beats back from now reaches a genuine beginning. On the infinite reading there is no first beat. The mesh has always been growing. Yet even then any self sees a finite past horizon, because only finitely many shells of docks lie inside its light cone. The signal from older shells has not had time to arrive.

\begin{proposition}[Finite age and horizon on both readings]
Either reading of birth yields a finite observable age and a finite past horizon for any self. On the finite reading this is so because the mesh itself is finite at every beat. On the infinite reading it is so because a self's light cone contains only finitely many shells. The two readings differ only in whether the very first shell exists.
\end{proposition}

The growth pays off either way. The only thing that differs is whether the first shell is real or whether the regress runs back without end. The model leaves that single question to the reader. What makes both readings work is that the mesh is finite at every beat, however far back the beats run. That finiteness is the load-bearing fact, and it is owned by the section on finiteness and growth.

\subsection{Unification is geometry}

The radial direction of the bulk, from the trunk outward toward the cusp, is the renormalization scale. This is the holographic reading of the growing mesh. Coarse-graining in physics, the act of looking at things at a larger scale and washing out the fine detail, corresponds to moving along the radial direction of the mesh. The deeper toward the trunk, the coarser the description.

\begin{principle}[The radial direction is the scale]
The bulk's radial direction is the renormalization scale. Moving inward toward the trunk is coarse-graining. Moving outward toward the cusp is refining. The running of the couplings with energy is this radial coarse-graining itself.
\end{principle}

From this one identification, force unification becomes a statement about shape. Deep in the bulk, near the trunk, the forces are one. Out on the flat cusp, near the twigs, they split into the separate forces we measure at everyday scales. The running of the coupling constants, the way the strengths of the forces drift as you change the energy of observation, is nothing but the radial coarse-graining of the mesh read as a change of scale.

\begin{result}[Grand unification is geometric]
Grand unification is the geometry of the mesh seen at the unification scale. The forces are one near the trunk and split toward the cusp. The unification is not an extra symmetry imposed by hand. It is a shape of the same growing mesh.
\end{result}

So the grand unification of the forces is not a separate symmetry grafted onto the model. It is the mesh seen from a particular radial depth. The same tree, read at the trunk, shows one force. Read at the twigs, it shows the several forces of the world.

\subsection{A computable cosmos}

The base is discrete and local. The mesh is a finite weave of docks at any beat, each dock carries 24 sites, each site holds one tone in $\{-1, 0, +1\}$, and the knit acts only on a dock and its immediate neighbours. There is no continuum anywhere in the base, and no value that is not one of three.

\begin{theorem}[Every finite region is computable]
Because the base is discrete and local, every finite region of the mesh over any finite span of beats is a finite computation. The universe is therefore simulable in principle, with no actual continuum and no uncomputable real numbers at the base.
\end{theorem}

This holds on both readings of birth. Whether the mesh is finite-growing or infinite-eternal, it is discrete at the base. The discreteness is the commitment that does the work here, and it is what makes the cosmos computable. The continuous structures that appear later in the paper, the smooth fields and the Lie groups, are emergent and coarse-grained, never present in the base. The careful version of the computability claim is carried by the section on computability and the refining universe.

\subsection{Consciousness as the substrate}

The base is felt. This is the one posit beyond the geometry and the knit. The vibe, the substance carried in each site and valued by the tone, is felt from the inside. A tone of $-1$ is pain, a tone of $0$ is peace, a tone of $+1$ is pleasure. These are not labels attached to a dead value. They are what the value is.

\begin{axiom}[The base is experiential]
The vibe is felt at the base. The single quale is one tone in $\{-1, 0, +1\}$, read as pain, peace, and pleasure. A dock of 24 sites is a hyper-color, the local texture of feeling. The substrate is experiential from the inside.
\end{axiom}

From this one posit the deepest implication follows directly. The universe is, at bottom, a structure of consciousness. Matter, space, gravity, and the forces are its shapes, the patterns that the felt substance falls into as the knit runs and the wake grows.

\begin{principle}[The carrying image, taken literally]
The universe is a mind, singing itself, burning as a fire, growing as a tree.
\end{principle}

The posit adds nothing to the base. The eight atoms (mesh, dock, site, tone, lean, knit, wake, beat) are unchanged by it. It only says that the substance those atoms are made of is felt from the inside rather than being a bare number. The geometry, the knit, and the growth are exactly as before. The single addition is that the value in each site is a feeling.

\subsection{Why it matters}

The implications are the payoff, not the trim. Each one is a place where several separate facts of physics collapse into a single feature of the base.

\begin{table}[ht]
\centering
\begin{tabular}{@{}ll@{}}
\toprule
implication & what collapses into one \\
\midrule
expansion and time are unified & flow of time, direction of time, and growth of space \\
a beginning and a horizon fall out & finite age and finite horizon, on either reading of birth \\
unification becomes geometry & the running couplings and the radial depth of the mesh \\
the cosmos is computable & no continuum, no uncomputable reals, simulable in principle \\
the base is experiential & matter, space, and forces are shapes of one felt substance \\
\bottomrule
\end{tabular}
\caption{The five implications and the separate facts each one unifies.}
\label{tab:implications}
\end{table}

The predictions of the model make it falsifiable, and the open problems say honestly what is unfinished. The implications say what is at stake if the model is right. The last of them, the experiential base, is the one that gives the model its felt weight. A computable, growing, unifying cosmos is already a strong claim. That the substance doing all of it is felt from the inside is what makes it a claim about the world we actually live in, and not only about a clever mechanism.

\begin{result}[What is at stake]
If the model holds, then time and expansion are one growth, a finite age and horizon come for free on either reading of birth, force unification is the geometry of the mesh, the cosmos is computable to its base, and the substance of all of it is consciousness. These are the payoff of the eight atoms, not additions to them.
\end{result}
